\section{Inżynieria silnika \amor{}}
\label{app:engineering}

Ten aneks nie jest dodatkowym wynikiem ekonomicznym. Jego rola jest inna: pokazuje, jakiego rodzaju infrastruktura obliczeniowa stoi pod eksperymentem BDP. Ma to znaczenie dla oceny wiarygodności pracy, ponieważ wynik nie pochodzi z arkusza kalkulacyjnego ani z jednorazowego skryptu, lecz z produkcyjnego silnika SFC-ABM rozwijanego jako projekt open-source. W bieżącym checkout repozytorium \amor{} warstwa Scala obejmuje 384 pliki i około 62 tys. linii kodu w \path{src/main/scala} oraz \path{src/test/scala}; z tego około 34 tys. linii przypada na kod silnika, a około 29 tys. na testy. Liczby te nie są argumentem z objętości, lecz wskazują skalę artefaktu, którego zachowanie jest następnie ograniczane przez typy, ledger, testy i dokumentację walidacyjną.

\subsection{Typy ekonomiczne i arytmetyka stałoprzecinkowa}

Silnik nie operuje na bezpostaciowych liczbach zmiennoprzecinkowych. Podstawowe wielkości ekonomiczne są reprezentowane jako typy stałoprzecinkowe oparte na \texttt{Long} ze skalą \(10^4\). W kodzie występują m.in. typy \texttt{PLN}, \texttt{Rate}, \texttt{Share}, \texttt{Scalar}, \texttt{Multiplier}, \texttt{Coefficient}, \texttt{PriceIndex} i \texttt{ExchangeRate}. Są one zdefiniowane jako typy \texttt{opaque}, co oznacza, że kompilator Scali nie pozwala traktować ich jako tej samej liczby. Przykładowo dodanie kwoty pieniężnej do stopy procentowej nie jest błędem runtime, lecz błędem typów jeszcze przed uruchomieniem programu.

Ten wybór projektowy ma dwa skutki. Po pierwsze, ogranicza klasę błędów jednostek, które w modelach makroekonomicznych są szczególnie trudne do wykrycia po fakcie. Po drugie, unika dryfu numerycznego typowego dla arytmetyki \texttt{Double}, który mógłby maskować lub generować pozorne naruszenia identyczności SFC. Tam, gdzie operacje mogą przekroczyć zakres, stosowane są jawne kontrakty arytmetyczne, np. \texttt{multiplyExact}, \texttt{addExact}, sprawdzane dzielenie i wersje z przejściem przez \texttt{BigInt} w ścieżkach wymagających ochrony przed przepełnieniem.

\subsection{Ledger i formalna weryfikacja}

\amor{} korzysta z osobnego modułu \ledger{}, utrzymywanego jako submoduł w \path{modules/ledger}. Ten moduł pełni funkcję jądra przepływów księgowych: agent lub mechanizm ekonomiczny może generować przepływ pieniężny, ale wykonanie przepływu przechodzi przez warstwę, której zadaniem jest utrzymanie zasady podwójnego zapisu.

Ważne jest precyzyjne rozróżnienie poziomów zaufania. Formalnie weryfikowany jest model referencyjny w \path{src/main/scala-stainless/Verified.scala}, sprawdzany przez Stainless oraz solver Z3. Dokumentacja \path{modules/ledger/docs/verification.md} wskazuje następujące własności tego modelu:

\begin{table}[htbp]
\centering
\scriptsize
\begin{tabularx}{\textwidth}{@{}L{0.30\textwidth}X@{}}
\toprule
\textbf{Własność} & \textbf{Znaczenie dla eksperymentu SFC-ABM} \\
\midrule
Konserwacja przepływu & suma sald dwóch rachunków dotkniętych przepływem pozostaje niezmieniona po jego wykonaniu \\
Warunek ramowy & rachunki nieuczestniczące w przepływie pozostają bez zmian \\
Aplikacja sekwencyjna & lista przepływów zachowuje własności konserwacji przy iteracyjnym wykonaniu \\
Dokładność dystrybucji & dystrybucja z resztą dokładnie sumuje się do kwoty źródłowej, bez utraty jednostek pieniężnych \\
Komutatywność przepływów rozłącznych & przepływy dotyczące rozłącznych rachunków dają ten sam wynik niezależnie od kolejności \\
Refinement runtime & model \texttt{Map[Int, Long]} jest powiązany z referencją \texttt{BigInt} pod jawnymi warunkami braku przepełnienia \\
\bottomrule
\end{tabularx}
\caption{Wybrane własności referencyjnego modelu ledgerowego weryfikowane przez Stainless/Z3.}
\label{tab:ledger-verification}
\end{table}

Produkcja nie jest w tekście przedstawiana jako w pełni formalnie udowodniona. To byłoby zbyt mocne twierdzenie. Produkcyjny interpreter, szybka ścieżka imperatywna i dystrybucja są testowane przeciwko współdzielonym modelom referencyjnym oraz kontraktom wykonania. Łańcuch zaufania wygląda zatem następująco: Stainless/Z3 dowodzi własności modelu referencyjnego; testy mostkują model referencyjny z czystym interpreterem produkcyjnym; testy równoważności porównują ścieżkę imperatywną z referencją; runtime wymusza kontrakty kształtu batchy, indeksów, nieujemnych przepływów i zakresów \texttt{Long}. Jest to celowo konserwatywna architektura: dowód formalny ma jasno określone granice, a wszystko poza nimi jest objęte testami i kontraktami.

\subsection{Warstwa agentów i mechanizmy gospodarki}

Warstwa agentowa \amor{} jest zapisana jako zbiór modułów z jawnym stanem i czystymi funkcjami przejścia. W katalogu \path{src/main/scala/com/boombustgroup/amorfati/agents} znajdują się m.in. moduły gospodarstw domowych, firm, sektora bankowego, NBP, ZUS/NFZ/PPK, JST, ubezpieczeń, NBFI, migracji, funduszy celowych, sektora quasi-fiskalnego oraz mechanizmów takich jak \texttt{EclStaging}, \texttt{DepositMobility} i \texttt{InterbankContagion}. Z punktu widzenia tej pracy ważne są zwłaszcza trzy elementy.

Po pierwsze, firmy nie są jedną reprezentatywną firmą. Mają stan technologiczny, kapitał, zapasy, zatrudnienie, gotówkę, dług i powody bankructwa. W kodzie istnieje m.in. przyczyna \texttt{AiDebtTrap}, czyli przypadek, w którym presja inwestycyjna i finansowanie technologii mogą doprowadzić firmę do bariery bilansowej. Po drugie, sektor bankowy nie jest pojedynczym agregatem. Model operuje na sześciu archetypach banków komercyjnych oraz agregacie \texttt{Others}; banki mają kapitał, NPL, LCR/NSFR, staging ECL zgodny z logiką IFRS 9 oraz możliwość upadłości. Po trzecie, ryzyko systemowe nie jest dopisane wyłącznie jako spread. Mechanizmy \texttt{InterbankContagion} i \texttt{DepositMobility} modelują straty kontrahentów, hoarding płynności oraz przepływ depozytów między bankami.

Ta architektura jest ważna dla interpretacji BDP. Jeżeli transfer zmienia płacę rezerwową, popyt, deficyt, koszt kapitału i kurs walutowy, to odpowiedź firm nie jest jedną funkcją agregatową. Jest wynikiem lokalnych decyzji, ograniczeń bilansowych, dostępu do finansowania i stanu technologicznego. Dzięki temu wynik nieliniowy nie musi być założony z góry. Może wyłonić się z interakcji kanałów.

\subsection{Powierzchnie walidacji i prowencja kalibracji}

\amor{} rozdziela co najmniej trzy powierzchnie kontroli. Pierwsza to walidacja księgowa: czy przepływy i zasoby zachowują identyczności SFC. Druga to artefakty macierzy SFC: Balance Sheet Matrix, Transactions Flow Matrix, stock-flow reconciliation oraz mapping z wierszy symbolicznych do mechanizmów runtime. Trzecia to walidacja empiryczna i kalibracyjna: czy wartości startowe, parametry i ścieżki modelu mają jawne źródła, statusy i ograniczenia.

Rejestr kalibracji \path{docs/calibration-register.md} jest generowany z typowanego rejestru prowencji. Nie jest listą luźnych przypisów, lecz tabelą parametrów z wartością, jednostką, właścicielem implementacyjnym, źródłem, transformacją i statusem. W bieżącej wersji rejestr wskazuje następujące liczby statusów:

\begin{table}[htbp]
\centering
\scriptsize
\begin{tabular}{@{}lr@{}}
\toprule
\textbf{Status prowencji} & \textbf{Liczba parametrów} \\
\midrule
\texttt{EMPIRICAL} & 36 \\
\texttt{EMPIRICAL\_TRANSFORMED} & 17 \\
\texttt{CODE\_NOTE\_EMPIRICAL} & 61 \\
\texttt{ASSUMED} & 31 \\
\texttt{TUNED\_NEEDS\_VALIDATION} & 85 \\
\texttt{POLICY\_SCENARIO} & 7 \\
\texttt{PLACEHOLDER} & 1 \\
\texttt{UNKNOWN\_SOURCE} & 0 \\
\bottomrule
\end{tabular}
\caption{Statusy prowencji kalibracji w produkcyjnym baseline \amor{} według rejestru kalibracji.}
\label{tab:calibration-provenance}
\end{table}

Takie oznaczenia są istotne metodologicznie. Model nie ukrywa, które parametry są empiryczne, które są transformacją danych instytucjonalnych, które są założeniem strukturalnym, a które są współczynnikami dostrojonymi i wymagającymi dalszej walidacji. To wzmacnia falsyfikowalność eksperymentu BDP: zewnętrzny badacz może nie tylko powtórzyć seedy, ale także wskazać, które parametry powinny być przedmiotem alternatywnej kalibracji lub analizy wrażliwości.

W tym sensie \amor{} jest nie tylko modelem ekonomicznym, ale także infrastrukturą replikowalnego sporu o mechanizmy. Twierdzenie „BDP jest warunkowym katalizatorem automatyzacji” nie jest zakończeniem analizy. Jest wynikiem jednego jawnego patcha, jednej jawnej siatki parametrów i jednego zestawu seedów na tle większego silnika, który można sklonować, zmodyfikować i uruchomić ponownie.
